\section{Place in the construction}

The longer physical blocking step is now defined, including the
source chart, its centered Gaussian, its gauge insertion and its
small-field cutoff. The pre-gauge representative of Notebook 25
belongs to this actual composite constraint. Choosing $R$ after its
primitive-base constants have been fixed can therefore provide a
strict amplitude margin for an inherited field on a fixed analytic
ball. This is the role of the new macrostep.

The next analytic task is to control the full inherited insertion
as strict joint activities in the coarse background and retained
fluctuations. A local-core expansion can keep its nearby field
dependence unexpanded and spend Taylor estimates only on its distant
tail. The precision and purely classical interactions can be
expanded directly in the coarse background before Gaussian
integration, using their existing finite classical support reserves.
Both operations require their explicit norm budgets; the pointwise
continuation proved here is not a substitute.

Even a common analytic domain and a convergent one-step logarithm
would leave a noncontracting activity-amplitude recurrence.
Extracting and controlling the vacuum and marginal gauge-action
terms is required for an arbitrarily long quantum iteration.
Large-field complements, mixed histories, observable estimates and
continuum reconstruction remain parts of the existence argument.

\begin{thebibliography}{9}
\bibitem{balaban109cm}
T.~Ba\l aban,
\emph{Renormalization group approach to lattice gauge field
theories. I},
Commun. Math. Phys. \textbf{109} (1987), 249--301.
\href{https://scispace.com/pdf/renormalization-group-approach-to-lattice-gauge-field-3ygai0vxt1.pdf}{Original article copy}.
The local ordinary average and contour mean are taken from
pp.~252--254; the composite-step arguments are developed here.
\bibitem{priorcm}
\emph{Toward a four-dimensional Yang--Mills construction},
research Notebooks 05--26 in this workspace.
The classical minimum and gauge identities, flat quadratic
locality, strict classical support families and the supported
pre-gauge representative are the previously established inputs.
\bibitem{sourcecm}
E.~Sussman and collaborators,
\href{https://github.com/ethanwsussman/Yang--Mills-LLM-experiment}{Yang--Mills-LLM-experiment},
deposited raw accounts at commit
\texttt{e392ef45013a581664bab76b90867170b30fe10e}.
\end{thebibliography}
\end{document}
